\documentclass[11pt]{article}
\usepackage[margin=0.85in]{geometry}
\usepackage{amsmath,amssymb}
\usepackage{booktabs}
\usepackage{array}
\usepackage{xcolor}
\usepackage{titlesec}
\usepackage[hidelinks]{hyperref}
\definecolor{good}{RGB}{20,120,60}
\definecolor{bad}{RGB}{180,40,40}
\definecolor{hdr}{RGB}{25,25,25}
\titlespacing*{\section}{0pt}{1.5ex}{0.6ex}
\titleformat{\section}{\large\bfseries\color{hdr}}{}{0pt}{}
\setlength{\parindent}{0pt}\setlength{\parskip}{0.7ex}
\newcommand{\g}[1]{\textcolor{good}{\textbf{#1}}}
\pagestyle{empty}

\begin{document}

\begin{center}
{\LARGE\bfseries When does the plan score matter? OGBench Cube}\\[3pt]
{\large A companion to the TwoRoom study: the same plan scores on a world model whose latent is \emph{already adequate}}\\[4pt]
{\small OGBench \textsf{cube-single} $\cdot$ \textsf{LeWM-cube}, \textsf{DINO-WM-cube} $\cdot$ goal-conditioned CEM-MPC}
\end{center}
\vspace{0.4em}

\section{Setup}
We repeat the TwoRoom plan-score comparison on OGBench \textbf{cube-single} (a 7-DoF arm must move a cube to a
target). We use the benchmark's own protocol: goal-conditioned MPC evaluated by replaying expert episodes and
setting the goal to the state \emph{$k$ steps ahead} (\emph{goal-offset} $k$), success $=$ cube within the goal
radius. Two frozen world models: \textbf{LeWM-cube} (pretrained) and \textbf{DINO-WM-cube} (a predictor we
trained; see caveat). Plan scores: \emph{latent} ($\lVert\hat z-z_g\rVert^2$, the default), \emph{TRM}
(regression metric), and \emph{TD value} (offline goal-conditioned temporal-difference, quasimetric head). All
numbers are \% of 40 goals reached; reference points: \textbf{expert $\approx$99\%}, published baselines
\textbf{LeWM 74 / DINO 86}, perfect planner 100. (Sanity: on 8 short goals every method scores 100\%.)

\section{1. Plan-score matrix}
\begin{center}\small
\renewcommand{\arraystretch}{1.25}
\begin{tabular}{ll ccc}
\toprule
\textbf{World model} & \textbf{goals} & \textbf{latent} & \textbf{TRM} & \textbf{TD value} \\
\midrule
\textsf{LeWM-cube} & easy (offset 25) & 80.0 & 82.5 & \g{85.0} \\
\textsf{LeWM-cube} & hard (offset 50) & 70.0 & 72.5 & 72.5 \\
\addlinespace[2pt]
\textsf{DINO-cube} & easy (offset 25) & 82.5 & 80.0 & 82.5 \\
\textsf{DINO-cube} & hard (offset 50) & 72.5 & \g{75.0} & 72.5 \\
\bottomrule
\end{tabular}
\end{center}
All four conditions sit in a tight band; the largest spread (LeWM easy: latent $80\!\to\!$ TD $85$, i.e. $+2$
successes in 40) is within $n{=}40$ sampling noise ($\pm$$\sim$7\%). Learned metrics neither clearly help nor hurt.

\section{2. TD-settings sweep (offset 25)}
\begin{center}\small
\renewcommand{\arraystretch}{1.2}
\begin{tabular}{l cccc c ccc}
\toprule
 & \multicolumn{4}{c}{expectile $\tau$ (at $n{=}20$)} & & \multicolumn{3}{c}{$n$-step (at $\tau{=}0.1$)} \\
\cmidrule(lr){2-5}\cmidrule(lr){7-9}
 & 0.1 & 0.3 & 0.5 & 0.7 & & 1 & 5 & 20 \\
\midrule
\textsf{LeWM} & 85 & 85 & 80 & 80 & & 85 & 85 & 85 \\
\textsf{DINO} & 82.5 & 85 & 87.5 & 87.5 & & 82.5 & 85 & 82.5 \\
\bottomrule
\end{tabular}
\end{center}
TD hyperparameters barely move success (all 80--87.5\%). The expectile --- the knob that \emph{doubled}
hard-goal success on TwoRoom --- is inert here; LeWM mildly prefers low $\tau$, DINO mildly high, both within noise.

\section{3. Planning-budget sweep (LeWM): does a weaker planner expose the plan score?}
\begin{center}\small
\renewcommand{\arraystretch}{1.2}
\begin{tabular}{l cc c cc}
\toprule
 & \multicolumn{2}{c}{easy (offset 25)} & & \multicolumn{2}{c}{hard (offset 50)} \\
\cmidrule(lr){2-3}\cmidrule(lr){5-6}
\textbf{CEM budget} & latent & TD value & & latent & TD value \\
\midrule
$10\times3$ ($\approx$30 evals) & 80.0 & 85.0 & & 72.5 & 72.5 \\
$30\times5$ ($\approx$150)      & 85.0 & 85.0 & & 70.0 & 70.0 \\
$60\times8$ ($\approx$480)      & 82.5 & 82.5 & & 72.5 & 70.0 \\
$150\times15$ (full)            & 80.0 & 85.0 & & 70.0 & 72.5 \\
\bottomrule
\end{tabular}
\end{center}
\textbf{No.} Even a \emph{tiny} $10\times3$ planner already reaches $\sim$80--85\%, and the TD-minus-latent gap
stays $\sim$0--5 (noise) at every budget. On TwoRoom a weak planner exposed the broken latent; on cube the
latent is good enough that minimal search suffices, so the plan score never becomes the bottleneck.

\section{4. Offline vs.\ online TD (LeWM, offset 25)}
\begin{center}\small
\renewcommand{\arraystretch}{1.2}
\begin{tabular}{l ccccc}
\toprule
 & offline & online $\tau$0.1 & online $\tau$0.3 & online $\tau$0.5 & online $\tau$0.7 \\
\midrule
success\,\% & 85.0 & 80.0 & 82.5 & 82.5 & 82.5 \\
\bottomrule
\end{tabular}
\end{center}
Online TD (the offline value continued on $\sim$2{,}500 planner-collected latents) $\approx$ offline at every
expectile --- no gain. On TwoRoom online TD helped because the offline value was poor on the planner's
(broken-latent, hard cross-wall) distribution; here the offline value is already good, so self-practice adds nothing.

\section{5. Twocube (two-cube manipulation)}
Does a \emph{harder} task push the latent back into the regime where the plan score pays off? We repeat the
comparison on OGBench \textbf{cube-double} (both cubes must reach targets) with \textsf{DINO-WM-double} (a
predictor we trained --- there is no pretrained double-cube world model, and no \textsf{LeWM-double}).

\emph{A methodological note first.} Our initial run rolled the predictor $25$ steps open-loop and every plan
score returned an \emph{identical} $37.5\%$ --- a red flag. The cause: this simplified residual-MLP predictor
\textbf{diverges} beyond $\sim$5 steps (terminal $\lVert z\rVert$: $1.1 \to 5\!\times\!10^{5}$ at 10 steps
$\to$ non-finite at 25), so the costs were garbage and equal only by coincidence. Following DINO-WM's
short-horizon \emph{receding} MPC, we retrain and roll only \textbf{5 steps} (the predictor's stable range),
which gives the valid numbers below.
\begin{center}\small
\renewcommand{\arraystretch}{1.25}
\begin{tabular}{l cccc}
\toprule
 & latent & TRM & TD value & online TD \\
\midrule
\textsf{DINO-WM-double} ($n{=}80$) & 31.2 & 30.0 & 31.2 & 32.5 \\
\bottomrule
\end{tabular}
\end{center}
The scores are now \textbf{statistically indistinguishable} ($30$--$32.5\%$; spread $\ll$ the $\pm$$\sim$5.5\%
band at $n{=}80$, though the success \emph{sets} differ, so the plan score does change behaviour). Success caps
at $\sim$31\% --- far below single-cube's $\sim$80\% --- because the short-horizon, simplified double-cube WM is
\textbf{capability-limited}: it solves only $\sim$31\% of two-cube goals regardless of terminal cost, leaving no
room for the plan score to help. Twocube is thus \textbf{inconclusive} on the plan-score question; settling it
needs a stronger double-cube model (the paper's $H{=}3$ ViT predictor, or a trained \textsf{LeWM-double}).

\section{Bottom line}
Across five axes --- learned metric vs.\ latent, TD hyperparameters, goal difficulty, planning budget, and
offline-vs-online --- single-cube success never leaves the $\sim$70--87\% band. This is the clean
\textbf{inverse of TwoRoom}, where a broken latent ($16\%$) made the learned reachability metric decisive
($\to$$68\%$). On the harder \textbf{twocube} the scores again cluster ($\sim$31\%), but for a different
reason --- a capability-limited world model, not an adequate latent. The unifying principle: \textbf{the
learned plan score is decisive in exactly one regime --- a \emph{misleading} latent with an
\emph{otherwise-capable} model+planner (TwoRoom). When the latent is already good (cube-single) or the model
itself is the bottleneck (twocube), the plan score is inert.}

\vspace{0.6em}\hrule\vspace{0.3em}
{\footnotesize\color{gray}
\textbf{Details.} Eval $=$ \texttt{World.evaluate} goal-offset replay, num\_eval 40, success via env. Plan-score
matrix / sweeps at CEM $150\times15$ unless noted. TD value $=$ offline goal-conditioned temporal-difference,
quasimetric head, $n{=}20$, $\gamma{=}0.99$. \textbf{Caveat:} \textsf{DINO-WM-cube} is a quick $\sim$6k-step
predictor we trained (action\_block 1), not a tuned checkpoint --- its row is the weaker leg; \textsf{LeWM-cube}
is the pretrained model and is the more reliable line. Twocube: \textsf{DINO-WM-double} rolled $5$ steps
(receding), $n{=}80$; the predictor is stable only for $\le$5-step rollouts (the paper uses an $H{=}3$ ViT
predictor). $n{=}40/80 \Rightarrow$ noise $\pm$$\sim$7\%/$5.5\%$.}

\end{document}
